% =============================================================================
%  Atto 8 — Lessons, Recap & What's Next  (~4 min, 3 slide)
% =============================================================================

\section{Lessons, recap \& what's next}

% ---------------------------------------------------------------------------- %
\begin{frame}{Recap: the whole picture in one image}
  \centering
  \includegraphics[width=0.78\textwidth]{MiviaRoverLogicalView.png}\\[2pt]
  \footnotesize
  \emph{Same picture as the very first slide of Part~II — but now every box
  is a thing you have seen the inside of.}\\
  Sensing $\to$ perception $\to$ Localization $\to$ Nav2 $\to$
  \texttt{ros2\_control} $\to$ DBC $\to$ firmware $\to$ wheels $\to$ encoders
  $\to$ /odom $\to$ Nav2 again.
\end{frame}

% ---------------------------------------------------------------------------- %
\begin{frame}{Lessons learned}
  \begin{enumerate}\setlength\itemsep{2pt}
    \item \textbf{The hardest bugs live at the boundaries.}\;
          RT-safety in \texttt{read()}, watchdog timing during bring-up,
          serialisation order across the bus — every painful day was at an
          interface, never inside a single component.
    \item \textbf{``Contract-first'' pays back fast.}\;
          The DBC paid for itself the first time SW and FW were rewritten
          in parallel, with no coordination call.
    \item \textbf{\texttt{diff\_drive\_controller} is enough for our 4WD
          skid-steer},\;provided wheels are grouped two-by-two and slip is
          absorbed at the firmware. No custom controller needed.
    \item \textbf{A custom \texttt{SystemInterface} is worth the effort.}\;
          Off-the-shelf solutions (e.g.\ \texttt{topic\_based\_ros2\_control})
          would skip the boundary, but trade away determinism and RT-safety.
  \end{enumerate}
\end{frame}

% ---------------------------------------------------------------------------- %
\begin{frame}{What's next \& hand-off}
  \begin{columns}[T,onlytextwidth]
    \column{0.55\textwidth}
      Two upstream blocks were \emph{assumed available} throughout this talk:

      \begin{itemize}\setlength\itemsep{3pt}
        \item \alert{Localization} — a planar EKF fusing wheel odometry,
              ZED IMU and a UWB anchor mesh; UWB quality is mapped into a
              dynamic covariance to handle NLOS.
        \item \alert{Perception} — Patchwork++ ground/non-ground
              segmentation on the ZED~2i point cloud, feeding the local
              costmap; ZED object detection feeds higher-level modules.
      \end{itemize}

      \vspace{0.4em}
      \alert{Hand-off.}\;
      The \emph{UWB-based Localization} story — including the heteroscedastic
      EKF — is the subject of \textbf{Emilio Amato's seminar, right after this one}.

    \column{0.45\textwidth}
      \textbf{Where to look in the repo.}
      \begin{itemize}\setlength\itemsep{2pt}
        \item \texttt{src/mivia\_rover\_platform/} — SystemInterface,
              URDF, controllers YAML, launch.
        \item \texttt{src/mivia\_rover\_can/} — DBC-aware messages, parser
              and command nodes, vendored \texttt{ros2\_socketcan}.
        \item \texttt{mivia\_rover\_firmware/} — STM32CubeIDE project,
              \texttt{rover.dbc}, FreeRTOS tasks.
        \item \texttt{src/mivia\_rover\_bringup/} — orchestrator with
              per-module enable flags.
      \end{itemize}
  \end{columns}

  \vspace{0.8em}
  \centering
  \alert{\textit{Thank you — questions?}}
\end{frame}
